(a) After translating $P$ to the origin, the linear term of $f$ is $f_x(P)x+f_y(P)y$. It is nonzero exactly when the Jacobian has rank $1$, which is equivalent to nonsingularity.

(b) At the origin the lowest homogeneous terms in \exref{I.5.1} are respectively $x^2$, $xy$, $-y^2$, and $xy(x+y)$. The multiplicities are therefore $2,2,2,3$. At either additional singular point of (c) in characteristics $7$ and $13$, the quadratic term after translation is $-\frac98(x-\frac34)^2+2(y-b)^2$, where $b^2=-1/2$. Both coefficients are nonzero, so these points also have multiplicity $2$.
